%%%%%%%%%%%%%%%%%%%%%%%%%%%%%%%%%%%%%%%%%
% Beamer Presentation
% Acelerador de Multiplicação de Matrizes em FPGA
%%%%%%%%%%%%%%%%%%%%%%%%%%%%%%%%%%%%%%%%%

\documentclass{beamer}

\mode<presentation>{
\usetheme{Madrid}

\definecolor{customorange}{RGB}{255, 140, 0}
\definecolor{customgreen}{RGB}{200, 245, 210}
\setbeamercolor{structure}{fg=customorange}
\setbeamercolor{item}{fg=customorange}
}

\usepackage{graphicx}
\usepackage{booktabs}
\usepackage[table]{xcolor}
\usepackage[brazilian]{babel}
\usepackage[utf8]{inputenc}
\usepackage[T1]{fontenc}
\usepackage{lmodern}
\usepackage{amsmath,amssymb}
\usepackage{url}
\usepackage{tikz}
\usetikzlibrary{arrows.meta,positioning,fit,calc}

\usepackage[
    backend=biber,
    style=numeric-comp,
    sorting=none,
    autocite=superscript
]{biblatex}

\addbibresource{refs.bib}

% Figuras numeradas automaticamente no Beamer
\setbeamertemplate{caption}[numbered]

% Fonte menor para a bibliografia
\renewcommand*{\bibfont}{\tiny}

% Figura opcional: se o arquivo existir, inclui; se não existir, mostra uma caixa indicando o caminho
\newcommand{\optionalfigurecaption}[4]{%
\begin{figure}
\centering
\IfFileExists{#1}{%
    \includegraphics[width=#2\linewidth]{#1}%
}{%
    \fbox{\parbox[c][0.30\textheight][c]{0.85\linewidth}{\centering Insira aqui a figura: \texttt{\detokenize{#1}}}}%
}%
\caption{#3}
\label{#4}
\end{figure}
}

\newcommand{\optionaltikz}[2]{%
\IfFileExists{#1}{%
    \resizebox{#2\linewidth}{!}{\input{#1}}%
}{%
    \fbox{\parbox[c][0.20\textheight][c]{#2\linewidth}{\centering Insira aqui o TikZ: \texttt{\detokenize{#1}}}}%
}%
}

\title[Acelerador FPGA]{Acelerador Parametrizável para Multiplicação de Matrizes em FPGA}
\subtitle{Estudo de caso usando DE0-CV e comunicação UART}
\author{Fernando Schettini \\ \textit{fernandoschettini@outlook.com}}
\institute[]{Estudo de caso e metodologia para exploração arquitetural em FPGA}
\date{2026}

\begin{document}

% ---------------------------------------------------------
% CAPA
% ---------------------------------------------------------

\begin{frame}
\titlepage
\end{frame}

\begin{frame}
\frametitle{Introdução - Motivação}

\begin{itemize}
    \item FPGAs são alternativas para offloading em sistemas embarcados e de borda, mas o desenvolvimento de aceleradores ainda depende de fluxos, plataformas e integrações específicas\supercite{haris2021secda,haris2023secdatflite}.
    
    \item Este trabalho explora uma arquitetura simples, parametrizável e mensurável para aceleração em FPGA de pequeno porte, composta por comunicação host-FPGA, armazenamento externo, buffers internos, unidades MAC e contadores de desempenho.
\end{itemize}

\end{frame}

\begin{frame}
\frametitle{Introdução - Por que multiplicação de matrizes?}

\begin{itemize}
    \item A multiplicação de matrizes é um kernel recorrente em IA, processamento de sinais e computação científica; em deep learning, operações do tipo GEMM aparecem em camadas densas, convolucionais e recorrentes\supercite{nvidia2023gemm}.
    
    \item O uso de matrizes \(128 \times 128\) com operandos INT8 representa uma carga compatível com inferência embarcada, em que a quantização inteira reduz tráfego de memória, latência e consumo de energia\supercite{tensorflow2022quantization}.
\end{itemize}

\end{frame}

\begin{frame}
\frametitle{Introdução - Objetivo}

\begin{itemize}
    \item Desenvolver um acelerador parametrizável de multiplicação de matrizes em FPGA usando a placa DE0-CV.
    \item Usar o PC como host, enviando matrizes para a FPGA por UART.
    \item Usar a FPGA como coprocessador para executar o cálculo.
\end{itemize}

\vspace{0.20cm}

\begin{center}
\optionaltikz{figs/host_fpga_uart.tex}{0.82}
\end{center}

\end{frame}

% =========================================================
\section{Fundamentação Teórica}
% =========================================================

\begin{frame}
\frametitle{Fundamentação Teórica - Definição do Problema}

\optionalfigurecaption
{figs/mac_parallel.jpeg}
{0.55}
{Representação conceitual de unidades de multiplicação e acumulação.}
{fig:mac-parallel}

\begin{itemize}
    \item Para duas matrizes $A$ e $B$, a matriz resultado $C$ é calculada por:
\end{itemize}

\[
    C_{ij} = \sum_{k=0}^{N-1} A_{ik}B_{kj}
\]

\end{frame}

\begin{frame}
\frametitle{Fundamentação Teórica - LABs, ALMs e caminhos locais}

\begin{columns}
\column{0.48\textwidth}
\optionalfigurecaption
{figs/lab_fast_local.png}
{0.85}
{Estrutura de LAB no Cyclone V\supercite{altera2023cyclonev}.}
{fig:lab}

\column{0.48\textwidth}
\optionalfigurecaption
{figs/alm_block.png}
{0.85}
{Estrutura de ALM no Cyclone V\supercite{altera2023cyclonev}.}
{fig:alm}
\end{columns}

\vspace{0.1cm}

\begin{itemize}
    \item LABs agrupam ALMs e interconexões locais\supercite{altera2023cyclonev}.
    \item ALMs implementam LUTs, registradores, somadores e caminhos de carry\supercite{altera2023cyclonev}.
    \item \textit{Logic element} é uma unidade lógica composta por LUT de quatro entradas, registrador programável e cadeia de carry\supercite{intel2017logicelement}.
\end{itemize}

\end{frame}

\begin{frame}
\frametitle{Fundamentação Teórica - Interconexões no Cyclone V}

\optionalfigurecaption
{figs/cyclonev_interconnects.png}
{0.70}
{Interconexões locais, diretas, por linha e por coluna no Cyclone V\supercite{altera2023cyclonev}.}
{fig:cyclonev-interconnects}

\vspace{0.1cm}

\begin{itemize}
    \item As interconexões locais, diretas, por linha e por coluna afetam o roteamento do circuito\supercite{altera2023cyclonev}.
    \item O roteamento pode se tornar um limite antes mesmo de todos os recursos lógicos serem usados.
\end{itemize}

\end{frame}

\begin{frame}
\frametitle{Fundamentação Teórica - PLLs}

\begin{itemize}
    \item PLLs permitem gerar clocks internos a partir do clock base da placa.
    \item São usados para ajustar a frequência de operação da lógica do acelerador e da interface SDRAM.
\end{itemize}

\end{frame}

\begin{frame}
\frametitle{Fundamentação Teórica - DSPs}

\begin{itemize}
    \item DSPs são blocos especializados para operações aritméticas, principalmente multiplicações.
    \item Neste trabalho, eles são explorados por meio de unidades MAC paralelas.
\end{itemize}

\vspace{0.15cm}

\optionalfigurecaption
{figs/dsp.png}
{0.60}
{Diagrama interno de um bloco DSP de precisão variável no Cyclone V\protect\supercite{altera2023cyclonev}.}
{fig:dsp-block}

\end{frame}

\begin{frame}
\frametitle{Fundamentação Teórica - Modos de operação dos DSPs}

\begin{table}
\centering
\tiny
\caption{Modos de operação dos blocos DSP de precisão variável no Cyclone V\protect\supercite{altera2023cyclonev}.}
\label{tab:dsp-modes}
\resizebox{\textwidth}{!}{%
\begin{tabular}{llllll}
\toprule
\textbf{Recurso} & \textbf{Modo} & \textbf{Instâncias} & \textbf{Pre-adder} & \textbf{Coeficiente} & \textbf{Chainout} \\
\midrule
1 DSP & 9 $\times$ 9 independente & 3 & Não & Não & Não \\
1 DSP & 18 $\times$ 18 independente & 2 & Sim & Sim & Não \\
1 DSP & 18 $\times$ 19 independente & 2 & Sim & Sim & Não \\
1 DSP & 18 $\times$ 25 independente & 1 & Sim & Sim & Sim \\
1 DSP & 20 $\times$ 24 independente & 1 & Sim & Sim & Sim \\
1 DSP & 27 $\times$ 27 independente & 1 & Sim & Sim & Sim \\
1 DSP & Dois 18 $\times$ 19 em modo adder & 1 & Sim & Sim & Sim \\
1 DSP & 18 $\times$ 18 somado com entrada de 36 bits & 1 & Sim & Não & Sim \\
2 DSPs & Multiplicação complexa 18 $\times$ 19 & 1 & Não & Não & Não \\
\bottomrule
\end{tabular}
}
\end{table}

\end{frame}

\begin{frame}
\frametitle{Fundamentação Teórica - Memória}

\begin{itemize}
    \item Blocos M10K são memórias embarcadas internas da FPGA dentro do chip.
    \item SDRAM é memória externa à FPGA, usada como armazenamento principal de maior capacidade.
    \item A arquitetura usa SDRAM para matrizes completas e M10K como buffers locais de tiles.
\end{itemize}

\end{frame}

\begin{frame}
\frametitle{Fundamentação Teórica - Recursos da DE0-CV}

\begin{itemize}
    \item A placa DE0-CV tem FPGA \textbf{Cyclone V 5CEBA4F23C7N}\supercite{terasic2015de0cv}.
    \item \textbf{49 mil} elementos lógicos programáveis, \textbf{18.480} Adaptive Logic Modules e \textbf{3080 Kbits} de memória embarcada\supercite{terasic2015de0cv}.
    \item 66 DSPs\supercite{terasic2015de0cv}.
    \item 4 PLLs fracionários\supercite{terasic2015de0cv}.
    \item 64 MB de SDRAM externa, com barramento de dados de 16 bits\supercite{terasic2015de0cv}.
\end{itemize}

\end{frame}

% =========================================================
\section{Metodologia}
% =========================================================

\begin{frame}
\frametitle{Metodologia - Compute Unit}

\begin{itemize}
    \item Multiplicação densa de matrizes inteiras:
    \[
        C = A \times B
    \]
    \item Entradas \(A\) e \(B\) em \textbf{INT8}; acumulador e saída \(C\) em \textbf{INT32}.
    \item Kernel baseado em tiles e parametrizado por \texttt{TILE\_SIZE}, \texttt{NUM\_MACS}, \texttt{DATA\_WIDTH} e \texttt{ACC\_WIDTH}.
    \item Uso de MACs paralelas para explorar os DSPs da FPGA.
\end{itemize}

\end{frame}

\begin{frame}
\frametitle{Metodologia - Fluxo local de cálculo}

\begin{itemize}
    \item Carregar tile de \(A\).
    \item Carregar tile de \(B\).
    \item Acumular os produtos no tile de \(C\).
    \item Escrever o resultado parcial ou final.
\end{itemize}

\vspace{0.15cm}

\begin{columns}[T]
\column{0.48\textwidth}
\centering
\optionaltikz{figs/compute_unit_io.tex}{1.0}

\vspace{0.1cm}
{\scriptsize \textbf{(a)} Interface externa da \texttt{compute\_unit\_top}.}

\column{0.48\textwidth}
\centering
\optionaltikz{figs/compute_unit_inside.tex}{1.0}

\vspace{0.1cm}
{\scriptsize \textbf{(b)} Organização interna da \texttt{compute\_unit\_top}.}
\end{columns}

\end{frame}

\begin{frame}
\frametitle{Metodologia - Otimização da Compute Unit}

\begin{itemize}
    \item A primeira fase otimiza apenas a compute unit.
    \item Objetivo: encontrar a configuração com maior GOPS efetivo.
    \item Cada experimento pertence a uma tupla discreta:
    \[
        (\texttt{TILE\_SIZE},\ \texttt{NUM\_MACS},\ \texttt{MAC\_PIPELINE\_STAGES})
    \]
    \item Foram avaliados:
    \begin{itemize}
        \item \texttt{TILE\_SIZE} em \(\{2,4,8,16\}\);
        \item \texttt{NUM\_MACS} em \(\{1,2,4,8,16,32,64\}\), conforme o tamanho do tile;
        \item \texttt{MAC\_PIPELINE\_STAGES} em \(\{0,1,2\}\), nas configurações selecionadas.
    \end{itemize}
\end{itemize}

\end{frame}

\begin{frame}
\frametitle{Metodologia - Como o GOPS é calculado}

\begin{itemize}
    \item Para matriz quadrada \(N \times N\), usa-se a aproximação:
    \[
        ops_{approx} = 2N^3
    \]

    \item Para \(N=128\):
    \[
        ops_{approx}=2\times128^3=4.194.304\ \text{operações}
    \]

    \item Tempo de execução interno:
    \[
        exec\_time = \frac{exec\_cycles}{f_{max}}
    \]

    \item GOPS efetivo:
    \[
        GOPS_{eff}=\frac{ops_{approx}}{exec\_time\times10^9}
    \]
\end{itemize}

\end{frame}

\begin{frame}
\frametitle{Metodologia - Fluxo experimental}

\begin{itemize}
    \item Para cada configuração, o fluxo automatizado executa:
    \begin{itemize}
        \item simulação;
        \item compilação no Quartus;
        \item extração de Fmax;
        \item extração de ALMs, DSPs e memória;
        \item geração de CSV final para comparação.
    \end{itemize}
\end{itemize}

\end{frame}

\begin{frame}
\frametitle{Metodologia - Organização de Memória}

\begin{itemize}
    \item As matrizes completas são armazenadas na SDRAM externa.
    \item Os tiles locais são armazenados em memórias internas M10K.
    \item A SDRAM atua como memória principal.
    \item As M10K funcionam como buffers rápidos próximos ao compute core.
    \item Organização em \textit{row-major}:
    \[
        addr = base + row \times N + col
    \]
\end{itemize}

\end{frame}

\begin{frame}
\frametitle{Metodologia - Espaços de memória}

\begin{itemize}
    \item O espaço de endereçamento é dividido em regiões separadas para:
    \begin{itemize}
        \item matriz \(A\);
        \item matriz \(B\);
        \item matriz \(C\).
    \end{itemize}
    \item Essa separação simplifica o controle de leitura, escrita e validação dos resultados.
\end{itemize}

\end{frame}

\begin{frame}
\frametitle{Metodologia - Otimização de Memória}

\begin{itemize}
    \item Segunda fase do \texttt{parameter\_optimization}.
    \item Objetivo: avaliar o impacto da organização de memória no desempenho final.
    \item Cada experimento pertence a uma tupla discreta:
    \[
        (\texttt{MEMORY\_BURST\_LEN},\ \texttt{BUFFERING\_MODE},\ 
        \texttt{MEMORY\_BANKS\_A},\ \texttt{MEMORY\_BANKS\_B})
    \]
    \item A varredura avalia 9 configurações pré-definidas.
    \item O modo \texttt{single} foi testado com 1 banco por matriz; o modo \texttt{double} foi testado com 2 ou 4 bancos.
\end{itemize}

\end{frame}

\begin{frame}
\frametitle{Metodologia - Critérios de comparação}

\begin{itemize}
    \item As configurações são comparadas por:
    \begin{itemize}
        \item GOPS efetivo;
        \item Fmax;
        \item uso de ALMs;
        \item uso de DSPs;
        \item uso de memória interna;
        \item pressão de roteamento.
    \end{itemize}
\end{itemize}

\end{frame}

\begin{frame}
\frametitle{Metodologia - Otimizações Implementadas}

\begin{itemize}
    \item Parametrização da arquitetura por generics.
    \item Separação entre top-level, compute, memória, controle e UART.
    \item Uso de SDRAM externa para matrizes completas.
    \item Uso de M10K como buffers locais de tiles.
    \item Uso de FIFO na UART para reduzir perda de bytes.
    \item PLL para clock da SDRAM.
\end{itemize}

\end{frame}

\begin{frame}
\frametitle{Metodologia - Otimizações Implementadas}

\begin{itemize}
    \item Leitura de \(C\) por checksum para reduzir o gargalo da UART.
    \item Suporte a \textit{single buffering} e \textit{double buffering}.
    \item Scripts automáticos para sweeps e coleta de métricas.
    \item Geração de CSV, gráficos e relatórios comparativos.
\end{itemize}

\end{frame}

% =========================================================
\section{Resultados e Discussões}
% =========================================================

\begin{frame}
\frametitle{Resultados - Top 3 da Compute Unit}

\begin{table}
\centering
\tiny
\caption{Melhores configurações do sweep \texttt{compute\_unit\_only}, ordenadas por \(GOPS_{eff}\) aproximado.}
\resizebox{\textwidth}{!}{%
\begin{tabular}{ccccccccc}
\toprule
\textbf{Run} & \textbf{Tile} & \textbf{MACs} & \textbf{Pipe} & \textbf{Fmax} & \textbf{DSPs} & \textbf{ALMs} & \textbf{GOPS} & \textbf{Risco} \\
\midrule
run\_013 & 16 & 64 & 0 & 65,13 MHz & 64 & 8144 & \textbf{7,9629} & alto \\
run\_017 & 16 & 64 & 1 & 74,68 MHz & 64 & 8156 & 7,3705 & alto \\
run\_018 & 16 & 64 & 2 & 79,23 MHz & 64 & 8148 & 6,5558 & alto \\
\bottomrule
\end{tabular}
}
\end{table}

\vspace{0.1cm}
\begin{itemize}
    \item O melhor desempenho puro veio de \texttt{TILE\_SIZE=16} e \texttt{NUM\_MACS=64}.
    \item Essa configuração usa praticamente todos os DSPs disponíveis e apresenta maior pressão de roteamento.
\end{itemize}

\end{frame}

\begin{frame}
\frametitle{Resultados - Top 3 do acelerador completo}

\begin{table}
\centering
\tiny
\caption{Melhores configurações do sweep com SDRAM/top completo. A linha verde indica a configuração escolhida para teste na placa.}
\resizebox{\textwidth}{!}{%
\begin{tabular}{ccccccccc}
\toprule
\textbf{Run} & \textbf{Buffer} & \textbf{Burst} & \textbf{Bancos A} & \textbf{Bancos B} & \textbf{Fmax} & \textbf{ALMs} & \textbf{Mem. bits} & \textbf{GOPS} \\
\midrule
\rowcolor{customgreen}
run\_005 & double & 4  & 2 & 2 & 72,00 MHz & 3433 & 68608  & \textbf{0,1613} \\
run\_006 & double & 8  & 2 & 2 & 72,00 MHz & 3433 & 68608  & 0,1613 \\
run\_007 & double & 16 & 2 & 2 & 72,00 MHz & 3433 & 68608  & 0,1613 \\
\bottomrule
\end{tabular}
}
\end{table}

\vspace{0.1cm}
\begin{itemize}
    \item Configuração escolhida: \texttt{N=128}, \texttt{TILE\_SIZE=8}, \texttt{NUM\_MACS=32}, \texttt{double buffering}, \texttt{MEMORY\_BURST\_LEN=4}, 2 bancos para A e 2 bancos para B.
    \item O melhor compute puro não foi escolhido diretamente porque o teste final inclui controle, SDRAM e UART.
\end{itemize}

\end{frame}

\begin{frame}
\frametitle{Resultados - Sweep previsto versus placa}

\begin{table}
\centering
\small
\caption{Comparação entre o resultado previsto pelo sweep e a execução real da \texttt{run\_20260609\_143308}.}
\resizebox{\textwidth}{!}{%
\begin{tabular}{lccc}
\toprule
\textbf{Métrica} & \textbf{Sweep full/top} & \textbf{Medição interna na placa} & \textbf{Fim-a-fim com UART} \\
\midrule
GOPS & 0,1613 & 0,0569 & 0,00125 \\
Percentual do previsto & 100\% & 35,3\% & 0,77\% \\
Clock usado & 72 MHz & 50 MHz & 50 MHz + UART \\
Ciclos internos & 1.872.130 & 3.684.540 & -- \\
Tempo medido & 26,00 ms & 73,69 ms & 16,84 s \\
\bottomrule
\end{tabular}
}
\end{table}

\vspace{0.05cm}
\begin{itemize}
    \item A queda para 35,3\% do previsto vem de dois fatores: clock real menor e quase 2 vezes mais ciclos internos.
    \item A queda fim-a-fim é dominada pela UART, pois o tempo inclui carregamento e retorno dos dados pelo host.
\end{itemize}

\end{frame}

\begin{frame}
\frametitle{Resultados - De onde vem a diferença de GOPS?}

\begin{columns}[T]
\column{0.48\textwidth}
\begin{block}{Modelo do sweep}
\begin{itemize}
    \item \(f_{max}=72\) MHz.
    \item \(exec\_cycles=1.872.130\).
    \item \(GOPS_{eff}=0,1613\).
\end{itemize}
\end{block}

\column{0.48\textwidth}
\begin{block}{Execução na placa}
\begin{itemize}
    \item Clock real: 50 MHz.
    \item Ciclos reais: 3.684.540.
    \item \(GOPS_{interno}=0,0569\).
\end{itemize}
\end{block}
\end{columns}

\vspace{0.2cm}
\[
    \frac{50}{72} \times \frac{1.872.130}{3.684.540}
    \approx 0,353
\]

\begin{itemize}
    \item Portanto, o desempenho interno medido corresponde a cerca de 35\% do valor previsto no sweep.
\end{itemize}

\end{frame}

\begin{frame}
\frametitle{Resultados - Comparação contextual com Raspberry Pi}

\begin{table}
\centering
\tiny
\caption{Comparação de ordem de grandeza. A linha do Raspberry Pi é referência externa de HPL/OpenBLAS em FP, não uma comparação direta com INT8.}
\resizebox{\textwidth}{!}{%
\begin{tabular}{lcccl}
\toprule
\textbf{Sistema / métrica} & \textbf{Tipo} & \textbf{Valor} & \textbf{Contra fim-a-fim UART} & \textbf{Observação} \\
\midrule
Compute unit FPGA, melhor sweep & INT8/INT32 & 7,9629 GOPS & 6370x & Sem SDRAM/UART \\
Acelerador completo, sweep & INT8/INT32 & 0,1613 GOPS & 129x & Modelo top completo \\
Acelerador na placa, interno & INT8/INT32 & 0,0569 GOPS & 45,5x & Sem contar transferência total \\
Acelerador fim-a-fim UART & INT8/INT32 & 0,00125 GOPS & 1x & Inclui comunicação UART \\
Raspberry Pi 5, HPL/OpenBLAS & FP / HPL & 24--25 GFLOP/s & 19.200x--20.000x & Referência contextual\supercite{rpi5hpl} \\
\bottomrule
\end{tabular}
}
\end{table}

\vspace{0.05cm}
\begin{itemize}
    \item A Raspberry Pi 5 usa CPU Broadcom BCM2712 quad-core Arm Cortex-A76 a 2,4 GHz\supercite{raspberrypi5official}.
    \item A comparação mostra que o gargalo atual não é só a compute unit, mas o sistema completo com UART e validação.
\end{itemize}

\end{frame}

\begin{frame}
\frametitle{Resultados - Por que o GOPS real ficou abaixo do esperado?}

\small

\begin{itemize}
    \item O sweep estimou a melhor configuração full/top com:
    \[
        GOPS_{sweep} = 0{,}1613
    \]

    \item Na placa, o clock real foi menor:
    \[
        \frac{50}{72} = 0{,}694
    \]

    \item A execução real também gastou mais ciclos:
    \[
        \frac{1.872.130}{3.684.540} = 0{,}508
    \]

    \item A queda observada é explicada por:
    \[
        0{,}694 \times 0{,}508 = 0{,}353
    \]
\end{itemize}

\vspace{0.1cm}

\begin{center}
\begin{tabular}{lc}
\toprule
\textbf{Métrica} & \textbf{Valor} \\
\midrule
GOPS esperado no sweep & 0,1613 \\
GOPS interno medido & 0,0569 \\
Percentual atingido & 35,3\% \\
GOPS fim-a-fim com UART & 0,00125 \\
\bottomrule
\end{tabular}
\end{center}

\vspace{0.05cm}

{\scriptsize
A run ainda não deve ser tratada como resultado final publicável, pois falhou validação:
\texttt{validation\_passed = false}, com 5 erros de checksum.
}

\end{frame}

\begin{frame}
\frametitle{Discussão - Gargalos observados}

\begin{itemize}
    \item A UART não é adequada para alto throughput de dados.
    \item A leitura completa de \(C\) pela UART se torna gargalo em matrizes grandes.
    \item O uso de checksum reduz o tempo de validação no host.
    \item O gargalo final pode migrar entre compute, memória e comunicação, dependendo da configuração.
\end{itemize}

\end{frame}

% =========================================================
\section{Conclusões e Prospecções}
% =========================================================

\begin{frame}
\frametitle{Conclusões}

\begin{itemize}
    \item Foi desenvolvido um acelerador parametrizável para multiplicação de matrizes inteiras em FPGA.
    \item A arquitetura permite avaliar diferentes configurações de compute e memória.
    \item O fluxo experimental automatiza simulação, síntese, coleta de métricas e comparação.
    \item A separação modular facilita a evolução do projeto.
\end{itemize}

\end{frame}

\begin{frame}
\frametitle{Prospecções}

\begin{itemize}
    \item Trocar a UART por uma interface de maior throughput.
    \item Criar uma interface padronizada para kernels.
    \item Permitir trocar o kernel sem reescrever controle e memória.
    \item Investigar se o gargalo principal está na memória, no protocolo ou no compute.
    \item Evoluir o uso de \textit{double buffering}.
\end{itemize}

\end{frame}

\begin{frame}
\frametitle{Prospecções - Acelerador como serviço}

\begin{itemize}
    \item Suportar múltiplos computadores ou clientes acessando o acelerador.
    \item Adicionar escalonamento de requisições.
    \item Adicionar isolamento entre usuários.
    \item Evoluir de um acelerador dedicado para uma infraestrutura compartilhável.
\end{itemize}

\end{frame}

% =========================================================
\section{Referências}
% =========================================================

\begin{frame}[allowframebreaks]
\frametitle{Referências}
\nocite{*}
\printbibliography
\end{frame}

% ---------------------------------------------------------
% FINAL
% ---------------------------------------------------------

\begin{frame}
\Large{\centerline{Obrigado pela presença!}}
\Large{\centerline{Perguntas?}}

\vspace{0.5cm}

\centering
{\small Fernando Schettini \\ \color{blue}fernandoschettini@outlook.com}

\end{frame}

\end{document}
